\chapter{Introduction}

\section{Contexte et justification}

Le secteur immobilier occupe une place croissante dans l'économie sénégalaise, porté par
l'urbanisation rapide de l'agglomération dakaroise et le développement de pôles urbains
secondaires (Diamniadio, Thiès, Mbour). Pourtant, contrairement à des marchés plus matures où
des bases de données de transactions publiques (cadastre ouvert, historiques notariaux)
permettent une estimation relativement transparente des prix, le marché sénégalais souffre
d'un déficit important d'information centralisée : les prix affichés sur les plateformes
d'annonces sont négociables, très hétérogènes d'un quartier à l'autre, et aucune base de
données publique ne permet aujourd'hui d'objectiver la valeur d'un bien à partir de ses
caractéristiques.

Cette opacité pénalise autant les particuliers (acheteurs, vendeurs) que les professionnels
(agences immobilières, banques accordant des prêts hypothécaires) qui doivent s'appuyer sur
l'expertise de terrain ou sur la méthode dite « des comparables », intrinsèquement subjective
et chronophage. L'essor de l'apprentissage automatique (\emph{machine learning}) et la
multiplication des données disponibles en ligne (notamment via les plateformes d'annonces
immobilières) ouvrent la voie à des modèles d'estimation automatique de la valeur des biens,
communément appelés \emph{Automated Valuation Models} (AVM) dans la littérature.

C'est dans ce contexte que s'inscrit ce projet informatique : construire, à partir de données
réellement observées sur le marché sénégalais, un modèle capable d'estimer le prix d'un bien
résidentiel (maison ou appartement), tout en explorant systématiquement l'apport de
l'optimisation des hyperparamètres d'entraînement sur la qualité des prédictions, et en rendant
ce modèle accessible via une interface graphique interactive plutôt que réservé à un usage
technique.

\section{Objectifs de la recherche}

Ce projet poursuit trois objectifs principaux :

\begin{enumerate}
  \item \textbf{Constituer un jeu de données exploitable} sur les prix des biens immobiliers
  résidentiels au Sénégal, en l'absence de source publique existante, par collecte automatisée
  de données d'annonces réelles.
  \item \textbf{Comparer plusieurs familles de modèles} de prédiction du prix (régression
  linéaire, méthodes ensemblistes à base d'arbres, réseau de neurones), en optimisant
  systématiquement leurs hyperparamètres d'entraînement par recherche aléatoire avec validation
  croisée, afin d'identifier l'approche la plus pertinente au regard du volume de données
  disponible.
  \item \textbf{Développer une interface graphique interactive} permettant à un utilisateur non
  technique d'obtenir une estimation de prix à partir des caractéristiques d'un bien, et
  d'explorer visuellement les données et les performances des modèles.
\end{enumerate}

\section{Structure du projet informatique}

Ce document est organisé en trois chapitres. Le \textbf{Chapitre~I} dresse un état de l'art du
domaine, à la fois sur les méthodes classiques d'évaluation immobilière et sur les travaux de
recherche appliquant l'apprentissage automatique à ce problème, en identifiant les zones encore
peu explorées, en particulier dans le contexte ouest-africain. Le \textbf{Chapitre~II} détaille
l'analyse des besoins, les sources de données mobilisées, le processus de collecte et de
pré-traitement, ainsi que la méthodologie de modélisation et d'évaluation retenue. Le
\textbf{Chapitre~III} présente l'implémentation réalisée : l'analyse exploratoire des données,
les résultats obtenus pour chaque modèle, leur interprétation et leur mise en perspective avec
la littérature, ainsi que les outils et l'architecture de la solution développée, incluant
l'interface graphique interactive. Une conclusion générale dresse le bilan du projet et propose
des perspectives d'amélioration.

\section{Conclusion}

Ce chapitre introductif a situé le projet dans son contexte : un marché immobilier sénégalais
en expansion mais peu documenté, pour lequel les techniques d'apprentissage automatique offrent
une opportunité concrète d'objectiver l'estimation des prix. Les trois objectifs retenus
(collecte de données, modélisation optimisée, restitution interactive) structurent l'ensemble
de la démarche présentée dans les chapitres suivants.
